% Generated by programs/20-build-analysis-report.R; do not edit.
\section{Empirical strategy}
Tables use effects per 10 matches; figures label their own units. Positive means more favorable to the target club. Main comparison intervals are pointwise 95\% season-block clustered intervals with small-sample t critical values. The two team rows from a match are not treated as independent games. Season blocks allow arbitrary dependence within a season but assume independence across seasons. Models compare recorded decisions, not correct and incorrect calls.
For Barcelona, L0 reports raw club indicators, L1 adds home, season, opponent and pre-match win-probability controls, and L2 adds splines in both teams' season-to-date possession and shots from strictly earlier matches. L3 instead adds same-match possession and shots to L1; L4 adds both teams' called fouls to L3 for card outcomes. The sample is held fixed across specifications within an outcome.
Barcelona's prior-match possession lies above the non-elite 95th percentile in 99.9\% of its eligible matches. Only 1.3\% lies inside Real Madrid's central 90\% range. The lagged-style models therefore depend on extrapolation, especially against non-elite clubs. L3 therefore uses same-match style instead of lagged style. It has more support but conditions on variables refereeing can change.
The Champions League models hold fixed venue, season, stage, opponent, Elo expected-score bins, possession bins and both teams' shots. The other-elite comparison excludes Barcelona, which has its own indicator. No domestic referee effects are transferred into this competition.
Pooling all eligible seasons estimates a persistent average gap rather than a different effect in every small season. It adds information without inventing matches or treating domestic and UEFA officials as interchangeable. The MDE is calibrated under a noncentral-t approximation at 80\% power and 5\% two-sided size, using the estimated clustered SE. Equivalence results use explicit illustrative bounds and 90\% intervals; failure to reject zero is not equivalence. Seasonal empirical-Bayes estimates retain the full covariance of jointly estimated season coefficients. Their intervals condition on estimated between-season heterogeneity and can understate uncertainty. Unpooled estimates, heterogeneity sensitivity, and leave-one-season-out fits remain available. Shrinkage intervals are not the headline evidence.
The before/after-2018 comparison is descriptive. The timing also coincides with VAR and changes in squads and performance. It is not an identified treatment effect and is not presented as a causal event study.
